\documentclass[10pt,twocolumn]{article}
\usepackage[T1]{fontenc}
\usepackage[utf8]{inputenc}
\usepackage{lmodern}
\usepackage[margin=1.8cm,top=2.4cm,bottom=2cm,columnsep=0.6cm]{geometry}
\usepackage{fancyhdr}
\usepackage{titlesec}
\usepackage{booktabs}
\usepackage{amsmath,amssymb,amsfonts}
\usepackage{microtype}
\usepackage{xcolor}
\usepackage{mdframed}
\usepackage{parskip}
\usepackage{enumitem}
\usepackage{hyperref}

\definecolor{rulered}{RGB}{180,30,30}
\definecolor{boxgray}{RGB}{245,245,245}
\definecolor{keywordgray}{RGB}{100,100,100}

\pagestyle{fancy}
\fancyhf{}
\fancyhead[L]{\textsc{\small Claude Mag}}
\fancyhead[C]{\small\textit{Machine Learning Research Digest}}
\fancyhead[R]{\small 2026-10-28}
\fancyfoot[C]{\small\thepage}
\renewcommand{\headrulewidth}{0.8pt}

\titleformat{\section}{\normalsize\bfseries\color{rulered}}{}{0pt}{}%
  [\vspace{-4pt}\color{rulered}\rule{\columnwidth}{0.4pt}]
\titlespacing{\section}{0pt}{8pt}{4pt}

\hypersetup{colorlinks=true,urlcolor=rulered,linkcolor=black}

\begin{document}
\setlength{\parindent}{0pt}
\setlength{\parskip}{4pt}

{\small\color{keywordgray}\textbf{Keywords:}
  omni-modal benchmarking \textbullet{}
  multimodal LLMs \textbullet{}
  ambiguity resolution \textbullet{}
  audio-visual integration}\par\vspace{4pt}

{\LARGE\bfseries\raggedright
  TRIAD: Benchmarking Omni-Modal Ambiguity in
  Multimodal Large Language Models}\par\vspace{4pt}

{\large\itshape\raggedright
  A New Benchmark Reveals That Even the Best Omni-Modal Systems
  Trail Humans by Over 20 Points When All Three Modalities Are
  Needed to Resolve a Scene}\par\vspace{6pt}

{\small\textbf{By} Zhaolu Kang, Yidi Wang, Yile Li, Siqi Zeng,
  Yingjie He, Richeng Xuan, Zhichao Hu}\par\vspace{2pt}

{\small\color{keywordgray}arXiv:2610.30841 \textbullet{}
  NeurIPS 2026 Datasets \& Benchmarks Track}
\par\vspace{2pt}\noindent{\color{rulered}\rule{\columnwidth}{0.6pt}}\vspace{6pt}

Multimodal language models that process text, images, and audio are
routinely evaluated on benchmarks where a single dominant modality
carries the answer---leaving intact an unmeasured gap between claimed
omni-modal capability and the genuine ability to fuse all three
modalities simultaneously. TRIAD closes this gap with 327 hand-crafted
queries spanning 86 scene groups and 18 ambiguity categories, each
designed so that the complete text--image--audio triplet yields a unique
answer while any single-modality subset leaves two or more options open.
Across 14 state-of-the-art omni-modal systems, the best model trails
humans by over 20 group-accuracy points, and removing answer options
drives model group accuracy to single digits.

\begin{mdframed}[backgroundcolor=boxgray,linecolor=rulered,linewidth=1pt,
    innerleftmargin=6pt,innerrightmargin=6pt,
    innertopmargin=6pt,innerbottommargin=6pt]
\textbf{\small AT A GLANCE}\par\vspace{4pt}
\begin{description}[leftmargin=0pt,labelsep=4pt,itemsep=2pt,parsep=0pt,
    font=\small\bfseries]
  \item[Problem:] Existing benchmarks do not require models to
    integrate all three modalities; a single modality usually suffices
    to answer correctly.
  \item[Why it matters:] Omni-modal systems are deployed as universal
    assistants; a model that silently ignores audio will fail in
    safety-critical settings where acoustic context is essential.
  \item[Method:] 327 text--image--audio triplets where the correct
    answer requires all three modalities, evaluated on item and scene-
    group accuracy with and without answer options.
  \item[Result:] Best model trails humans by $>$20 points on group
    accuracy; removing answer options drops model group accuracy to
    single digits.
  \item[Evidence:] 14 state-of-the-art omni-modal systems on 86 scene
    groups covering 18 ambiguity categories.
\end{description}
\end{mdframed}

\section{The Big Picture}

The ambition of omni-modal language models---processing text, image,
and audio within a single system---promises agents that perceive the
world as richly as humans do. Yet evaluating this ambition has relied
on benchmarks that are functionally bimodal: image-question answering
where text alone is sufficient, or audio tasks where the image adds
marginal signal. True omni-modal integration requires all three
modalities to be jointly necessary.

TRIAD operationalizes this requirement through the principle of
\emph{modal indispensability}: for each query, every modality is
necessary and no proper subset is sufficient. The dataset captures
everyday scenes where text, image, and audio together resolve what
neither pair can---e.g., hearing an alarm while seeing a clock while
reading a date.

\section{Why Existing Approaches Fall Short}

\textbf{Item-level evaluation is misleading.}
A model that achieves 70\% accuracy on individual items may still fail
every scene group, because group accuracy requires consistent
disambiguation across all items sharing a scene. TRIAD reports both
metrics; the gap between them exposes unreliable partial reasoning.

\textbf{Answer options inflate performance.}
Multiple-choice formats let models exploit statistical regularities in
the option set (e.g., always picking the longest option). TRIAD
evaluates both with and without answer options; performance drops to
single digits in the open-ended condition, confirming option exploitation.

\textbf{Audio cues are underweighted.}
Models trained on predominantly text--vision data default to text--image
reasoning even when the disambiguating signal is acoustic. TRIAD's
audio-indispensable categories expose this systematic underweighting.

\section{Core Mechanism}

TRIAD's benchmark construction follows three stages:

\begin{enumerate}[itemsep=2pt,parsep=0pt]
  \item \textbf{Scene collection:} Authors collect or synthesize
    naturalistic scenes representable in all three modalities with
    complementary but partially overlapping information.
  \item \textbf{Ambiguity verification:} Each triplet is manually
    reviewed to confirm that any two-modality combination leaves
    at least two consistent answer options.
  \item \textbf{Query authoring:} Questions are written to be
    answerable only with the full triplet; distractors are crafted
    to be plausible given incomplete modality access.
\end{enumerate}

Scene groups (86 in total) share a common real-world context and
each contain 3--5 queries, allowing the group accuracy metric to
test whether disambiguation generalizes within a scenario.

\section{Technical Formulation}

Let $\mathcal{T} = (t, v, a)$ denote a triplet of text $t$, image $v$,
and audio $a$, and let $\mathcal{A} = \{a_1, \ldots, a_K\}$ be the
answer set. A triplet is modally indispensable if:
\[
|\arg\max_k P(a_k \mid \mathcal{T})| = 1
\quad\text{but}\quad
|\arg\max_k P(a_k \mid \mathcal{T} \setminus \{m\})| > 1
\]
for every modality $m \in \{t, v, a\}$. This condition defines the 327
items in TRIAD; it is verified by human annotators, not automatically.

Group accuracy for scene group $g$ is:
\[
\text{GroupAcc}(g) = \mathbf{1}\!\left[\forall q \in g,\;
  \hat{a}(q) = a^*(q)\right]
\]
where $\hat{a}(q)$ is the model's predicted answer and $a^*(q)$ the
ground-truth. Mean group accuracy across all 86 groups is the primary
metric.

\section{Learning or Inference Procedure}

TRIAD is a static evaluation benchmark; no training is involved.
Models are prompted with the triplet in their native input format
(text + image + audio), asked to output either a multiple-choice
letter or a free-form answer, and scored automatically. Human
performance is estimated from a held-out annotator pool distinct
from the construction team.

\section{What the Guarantee Says}

The modal-indispensability design guarantee is:
\begin{itemize}[itemsep=2pt,parsep=0pt]
  \item \textbf{Soundness:} A model that answers correctly must have
    accessed all three modalities (no single-modality shortcut exists
    by construction).
  \item \textbf{Completeness:} Group accuracy measures consistent
    disambiguation within a scene, ruling out lucky individual item
    guesses.
\end{itemize}
These are guarantees about the benchmark's diagnostic validity, not
about any particular model's capability.

\section{Experimental Findings}

\begin{itemize}[itemsep=2pt,parsep=0pt]
  \item \textbf{Group accuracy gap:} The best of 14 evaluated omni-modal
    systems achieves group accuracy $>$20 points below human performance.
  \item \textbf{Item vs.\ group accuracy:} Models with high item accuracy
    (e.g., $>$65\%) still fail most scene groups, confirming inconsistency.
  \item \textbf{Open-ended collapse:} Removing answer options drops average
    model group accuracy to single digits; human group accuracy declines
    only modestly (from $\sim$85\% to $\sim$70\%).
  \item \textbf{Audio underuse:} On audio-indispensable items, models score
    near chance even in the multiple-choice condition, confirming that
    acoustic reasoning is the primary failure mode.
\end{itemize}

\section{Ablations and Interpretation}

\textbf{Modality ablation:}
Providing models with only text+image, text+audio, or image+audio gives
performance close to random, confirming benchmark validity.

\textbf{Model family comparison:}
Closed-source frontier models outperform open-source models by a
larger margin on TRIAD than on existing benchmarks, suggesting that
scale and instruction tuning improve omni-modal integration but do not
close the gap.

\textbf{Category breakdown:}
The hardest ambiguity categories involve temporal audio events (e.g.,
distinguishing two events from their sounds alone) and spatial
sound-source localization, pointing toward specific architectural
improvements.

\section*{Reference}

Zhaolu Kang, Yidi Wang, Yile Li, Siqi Zeng, Yingjie He, Richeng Xuan,
Zhichao Hu.
\textbf{TRIAD: Benchmarking Omni-Modal Ambiguity in Multimodal Large
Language Models}.
arXiv:2610.30841. NeurIPS 2026.
\url{https://arxiv.org/abs/2610.30841}

\end{document}
